\clearpage\thispagestyle{empty}
\pdfbookmark[0]{Quellen zu Pallas (29. Mai 1814)}{pallas1814may-sources}
\begin{center}\textbf{Quellen zu Pallas (29. Mai 1814)}\end{center}
\noindent Diese KI-Übertragung ins heutige Deutsch folgt dem Brief von Carl Friedrich Gauß mit der Überschrift Pallas in seinen \href{https://gdz.sub.uni-goettingen.de/id/PPN236007467}{\textit{Werken}, Band~VI} (Göttingen, 1874), S.~376--377. Gauß sandte ihn am 29. Mai 1814 aus Göttingen an Johann Elert Bode, den Herausgeber des \textit{Astronomischen Jahrbuchs}. Der Brief erschien im Jahrbuch für 1817 (Berlin, 1814), S.~\href{https://api.digitale-sammlungen.de/iiif/image/v2/bsb10538346_00220/full/1600,/0/default.jpg}{212}, \href{https://api.digitale-sammlungen.de/iiif/image/v2/bsb10538346_00221/full/1600,/0/default.jpg}{213}, \href{https://api.digitale-sammlungen.de/iiif/image/v2/bsb10538346_00222/full/1600,/0/default.jpg}{214}, \href{https://api.digitale-sammlungen.de/iiif/image/v2/bsb10538346_00223/full/1600,/0/default.jpg}{215}. Die Links führen zum digitalisierten Exemplar der Bayerischen Staatsbibliothek.
\par\smallskip\noindent Die Überschrift des Erstdrucks nennt die Beobachtungen der Juno und ihre neuen Bahnelemente, die nächste Opposition der Pallas und astronomische Nachrichten. Eine mit B. gezeichnete Anmerkung auf S.~212 datiert das frühere Schreiben zur Juno auf den 10. Dezember 1813. In der Werkausgabe steht der Brief unter Pallas. Die längere Überschrift und die Anmerkung fehlen dort. In beiden Ausgaben lautet der Name des Instrumentenbauers Frauenhofer. KI-Übertragung und KI-Satz.
\clearpage
